%=========== ΛΥΣΗ - 13ο ΘΕΜΑ Α ===========
\providecommand{\theoriaref}[3]{%
\begin{center}
\fcolorbox{maincolor!70!black}{maincolor!8}{%
\begin{minipage}{.88\linewidth}
\textcolor{maincolor!80!black}{\kerkissans{\textbf{#1~\ref{#2}}}}\\[1mm]
#3
\end{minipage}}
\end{center}}
\providecommand{\orismosref}[1]{%
\theoriaref{Ορισμός}{#1}{Η απάντηση δίνεται από τον αντίστοιχο ορισμό.}}
\providecommand{\apodeixiref}[1]{%
\theoriaref{Θεώρημα}{#1}{Η ζητούμενη απόδειξη δίνεται στην αντίστοιχη απόδειξη.}}

\begin{thema}{Α}
\begin{erwthma}
\item \apodeixiref{thewrhma:paragwgos_logarithmou}

\item \orismosref{orismos:kyrth_synarthsh}

\item \orismosref{orismos:paragwgisimh_se_diasthma}

\item Οι χαρακτηρισμοί των προτάσεων είναι:
\begin{alist}
\item \textbf{Σωστό}. Από το Θ.Μ.Τ. υπάρχει $\xi\in(a,\beta)$ τέτοιο ώστε
\[ f'(\xi)=\dfrac{f(\beta)-f(a)}{\beta-a}. \]
Όμως
\[ \dfrac{f(a)-f(\beta)}{a-\beta}=\dfrac{f(\beta)-f(a)}{\beta-a}. \]
\item \textbf{Σωστό}. Το ορισμένο ολοκλήρωμα είναι γραμμικό, άρα
\[ \dint_a^\beta (f(x)+g(x))\d x=\dint_a^\beta f(x)\d x+\dint_a^\beta g(x)\d x. \]
\item \textbf{Λάθος}. Η σχέση $f'(x_0)=0$ δεν αρκεί για τοπικό ακρότατο. Για παράδειγμα, για $f(x)=x^3$ έχουμε $f'(0)=0$, αλλά το $0$ δεν είναι θέση τοπικού ακρότατου.
\item \textbf{Σωστό}. Κάθε συνεχής συνάρτηση στο $\mathbb{R}$ έχει παράγουσα, για παράδειγμα τη συνάρτηση
\[ F(x)=\dint_0^x f(t)\d t. \]
\item \textbf{Λάθος}. Ίσες συναρτήσεις χρειάζεται να έχουν ίδιο πεδίο ορισμού και ίδιες τιμές, όχι απαραίτητα ίδιο τύπο. Για παράδειγμα, οι τύποι $f(x)=x^2$ και $g(x)=|x|^2$ ορίζουν την ίδια συνάρτηση στο $\mathbb{R}$.
\end{alist}

\item Σωστή απάντηση είναι η \textbf{β}. Σύμφωνα με το θεώρημα \eng{Bolzano}, αν $f$ είναι συνεχής στο $[a,\beta]$ και $f(a)\cdot f(\beta)<0$, τότε υπάρχει τουλάχιστον μία ρίζα στο $(a,\beta)$.
\end{erwthma}
\end{thema}
%=========== ΤΕΛΟΣ ΛΥΣΗΣ - 13ο ΘΕΜΑ Α ===========
